% !TeX root = ../main.tex

\chapter{结论与展望}

\section{主要结论}

针对四槽并联碱性电解水（AWE）系统的协调控制问题，
提出了基于扩散模型的安全预测控制框架，并完成了系统的理论研究与仿真验证。

\subsection{性能总结}

从控制精度、安全性与适用性三个角度，实验结果可归纳如下：

\begin{enumerate}
  \item \textbf{Diffusion TCN策略表现优异}：
  所提出的Diffusion TCN策略在功率跟踪精度上接近NMPC基线，
  功率跟踪RMSE为0.14 MW，温度调节MAE为2.0 K。
  同时，推理时间从NMPC的2--5秒降低至约50毫秒，实现了近两个数量级的计算加速，
  满足了实时控制的需求。

  \item \textbf{安全约束保证机制有效可靠}：
  引入投影算子与CBF安全滤波后，系统严格满足各项安全约束。
  HTO始终低于2\%的安全阈值，温度、电压与功率约束满足率均达到100\%，
  验证了两种安全增强机制的有效性。

  \item \textbf{生成模型适用于工业控制}：
  扩散模型等生成方法不仅可学习复杂控制策略，
  还可借助代价加权动作选择机制实现在线优化，
  展示了生成模型在工业过程控制中的应用潜力。
\end{enumerate}

\subsection{方法对比总结}

表~\ref{tab:final_comparison} 总结不同方法的特点和适用场景。

\begin{table}[!htbp]
  \centering
  \caption{控制方法综合对比}
  \label{tab:final_comparison}
  \begin{tabular}{@{}l c c c p{3.5cm}@{}}
    \toprule
    \textbf{方法} & \textbf{跟踪精度} & \textbf{计算效率} & \textbf{安全保证} & \textbf{适用场景} \\
    \midrule
    NMPC & 高 & 低 & 中 & 离线规划、慢速过程 \\
    \addlinespace
    Diffusion & 高 & 高 & 中 & 需要生成式动作选择的场景 \\
    \addlinespace
    Diffusion+安全投影 & 高 & 高 & 高 & 安全关键型实时控制 \\
    \addlinespace
    Diffusion+CBF & 高 & 高 & 高 & 需要理论安全保证的场景 \\
    \bottomrule
  \end{tabular}
\end{table}

由表~\ref{tab:final_comparison} 可知，NMPC适合对实时性要求不高的离线规划，
而Diffusion TCN在保持较高跟踪精度的同时，将推理时间降至约50毫秒；
叠加安全投影或CBF滤波后，均在毫秒级时间内实现了严格的安全保证，
适宜安全关键型实时控制场景。

\section{创新点}

本研究的主要创新点包括：

\begin{enumerate}
  \item \textbf{面向多槽AWE系统的Diffusion TCN预测控制策略}：
  针对多槽电解系统的功率分配与温度调节问题，提出基于Diffusion TCN的生成式预测控制策略，
  利用扩散模型的采样能力捕捉多种等价最优动作模式，将单次推理时间从NMPC的2--5秒降至约50毫秒。

  \item \textbf{非仿射CBF安全滤波与严格安全保证}：
  针对生成模型输出的尾部风险，建立了面向多槽AWE系统的非线性非仿射CBF框架，
  显式处理了电流-法拉第效率的非仿射耦合项，从理论上证明了安全集的前向不变性，
  从理论上保证了安全集的前向不变性。

  \item \textbf{安全投影算子与CBF滤波的对比验证}：
  实现了投影算子单步预测修正与CBF李导数理论保证两类并行安全机制，
  在正常工况与高温、高功率、低功率HTO累积等多类极端场景下完成了对比验证，
  量化了两者在约束满足率、计算开销和控制品质上的差异。
\end{enumerate}

\section{局限性与未来工作}

\subsection{研究局限性}

尽管研究在仿真层面验证了所提控制方法的有效性与实时性，但仍存在以下局限性：

\begin{enumerate}
  \item \textbf{仅限于仿真验证}：
  当前研究完全基于计算机仿真，未在真实电解系统上开展实验验证。
  仿真模型与实际系统之间存在差异，所提方法在实际应用中的有效性有待检验。

  \item \textbf{模型泛化能力有限}：
  训练得到的策略模型对未知工作条件和设备老化情况的泛化能力尚需进一步研究。

  \item \textbf{安全约束范围有限}：
  当前安全约束处理机制仅考虑部分关键安全约束，对复杂故障模式和极端工况的处理能力有限。
\end{enumerate}

\subsection{未来研究方向}

基于上述局限性，未来研究可从以下方向展开：

\begin{enumerate}
  \item \textbf{Sim-to-Real迁移}：
  开展数字孪生到真实电解槽的迁移研究，通过域随机化、系统辨识等方法
  缩小仿真与现实的差距，最终实现真实系统的部署验证。

  \item \textbf{在线自适应更新}：
  研究策略模型的在线自适应更新机制，使控制器能够应对设备老化、
  工况变化等时变特性，提高系统的长期运行稳定性。

  \item \textbf{多目标优化}：
  在功率跟踪和温度调节的基础上，进一步考虑氢气纯度、系统能效等多目标优化问题，
  实现经济效益与安全运行的综合优化。

  \item \textbf{多槽启停调度与运行状态管理}：
  当前研究假设四槽始终处于正常运行状态，未考虑部分槽体启停调度问题。
  实际工业运行中，电解槽可能经历冷启动、热启动、热待机与正常运行等多种状态，
  各状态下的热惯性、电化学响应速度与安全边界存在显著差异。
  未来可研究融合多槽启停决策与扩散模型预测控制的协同调度框架，
  在宽功率波动场景下根据负载水平动态决定各槽运行状态，
  兼顾系统响应速度、设备寿命与整体能效。

  \item \textbf{与其他控制方法的融合}：
  探索生成模型与传统MPC、自适应控制等方法的融合，
  结合各自优势构建高性能混合控制框架。
\end{enumerate}

\section{总结}

针对四槽并联AWE系统在波动性可再生能源接入背景下的协调控制难题，
本文建立了涵盖电化学、热力学与氢气传输的高保真耦合动力学模型，
设计了NMPC专家数据生成机制，提出了基于Diffusion TCN的生成式预测控制策略，
并引入投影算子与CBF安全滤波两类约束保证机制，
构建了完整的仿真-训练-验证实验平台。

实验结果表明，Diffusion TCN策略在保持接近NMPC基线控制精度的同时，
单次推理时间从2--5秒降至约50毫秒，实现了近两个数量级的计算加速。
更关键的是，叠加安全投影或CBF滤波后，系统在全部测试时段内严格满足温度、电压、功率与HTO约束，
未发生任何越界事件，验证了生成模型与安全机制融合的有效性。
特别地，CBF框架通过显式处理电流-法拉第效率的非仿射耦合，
从理论上证明了安全集的前向不变性。
上述成果为AWE系统与波动性可再生能源的协调运行提供了兼顾控制精度、计算效率与安全保证的工程化方案。
